\section{演習問題}
ここではテキストの内容を活かした効率的な解法を実際に行ってみよう.目的は読んだ内容が本当に成り立つかを
確認するためなので問題数は多く用意していない.読書は”\textcolor{red}{絶対に}”発展内容まで読んでいるので,
前提とした解説を行う.\\
\subsection{rank}
\subsubsection{問題}
\noindent
\fbox{1}\\
rank
$\begin{pmatrix}
    2 & 3 & 1 & 7 \\
    1 & 1 & 1 & 2 \\
    3 & 2 & 4 & 5
\end{pmatrix}$ \hspace{1cm}
rank$\begin{pmatrix}
    2 & 1 & 1 & 3 \\
    3 & 2 & 1 & 2 \\
    1 & 1 & 2 & 1 \\
    2 & 1 & 3 & a
\end{pmatrix}$\\
\begin{flushright}
    （2025線形代数学1a）
\end{flushright}
\fbox{2}\\
階数が1となるように,以下の行列の空いているところを埋めよ：
\[
A=
\begin{bmatrix}
    1 & 2 & 4 \\
    2 & & \\
    4 & & 
\end{bmatrix} \qquad
B=
\begin{bmatrix}
    & 9 &  \\
    1 & & \\
    2 & 6 & -3
\end{bmatrix}
\]
\begin{flushright}
    （「ストラング：線形代数イントロダクション」p161）
\end{flushright}
\subsubsection{解答}
\noindent
\fbox{1}\\
rank
$\begin{pmatrix}
    2 & 3 & 1 & 7 \\
    1 & 1 & 1 & 2 \\
    3 & 2 & 4 & 5
\end{pmatrix}\rightarrow$
rank
$\begin{pmatrix}
    0 & 1 & -1 & 3 \\
    1 & 1 & 1 & 2 \\
    0 & -1 & 1 & -1
\end{pmatrix}\rightarrow$rank
$\begin{pmatrix}
    1 & 0 & 2 & 1 \\
    0 & 1 & -1 & 1 \\
    0 & 0 & 0 & 2
\end{pmatrix}=3$\\
rank
$\begin{pmatrix}
    2 & 1 & 1 & 3 \\
    3 & 2 & 1 & 2 \\
    1 & 1 & 2 & 1 \\
    2 & 1 & 3 & a
\end{pmatrix}\rightarrow$rank
$\begin{pmatrix}
    2 & 1 & 1 & 3 \\
    1 & 1 & 0 & -1 \\
    1 & 1 & 2 & 1 \\
    0 & 0 & 2 & a-3
\end{pmatrix}\rightarrow$rank
$\begin{pmatrix}
    0 & -1 & 1 & 5 \\
    1 & 1 & 0 & -1 \\
    0 & 0 & 2 & 2 \\
    0 & 0 & 2 & a-3
\end{pmatrix}$\\ 
\begin{flushright}
$\rightarrow \begin{pmatrix}
    1 & 1 & 0 & -1 \\
    0 & 1 & -1 & -5 \\
    0 & 0 & 1 & 1 \\
    0 & 0 & 0 & a-5
\end{pmatrix}
\quad \therefore
\begin{cases*}
    3 \quad (a=5) \\
    4 \quad (a \neq 5)
\end{cases*}$
\end{flushright}
\fbox{2}\\
行列のサイズによらず,第1列にのみpivotがあれば良い.\\
Aは第2行,第3行が[1 2 4]の倍数になれば消える.\\
よって
\[A=
\begin{bmatrix}
    1 & 2 & 4 \\
    2 & 4 & 8 \\
    4 & 8 & 18
\end{bmatrix}\]
第2列の1を活かしたい.\\
\[B=
\begin{bmatrix}
    & 9 & \\
    1 & 3 & -3/2 \\
    2 & 6 & -3    
\end{bmatrix}\rightarrow
\begin{bmatrix}
    3 & 9 & -9/2 \\
    1 & 3 & -3/2 \\
    2 & 6 & -3  
\end{bmatrix}\]
\subsection{一次方程式の求解}
\subsubsection{問題}

\noindent
\fbox{1}
以下の空欄を適切な数式または数値で埋めよ.パタメーターを要する場合は,文字s,t,u,...をこの順に使用し,「標準的なパラメーターの置き方」に従うこと.\\
(1)連立一次方程式
$\begin{cases*}
    4x+8y=12\\
    7x+12y=21
\end{cases*}$の解は$x=$\fbox{\makebox[1.0cm][c]{1}},$y=$\fbox{\makebox[1.0cm][c]{2}}である.\\
(2)連立一次方程式
$\begin{cases*}
    3x_1+3_x2+x_3+10x_4=15 \\
    2x_1+3x_2+7x_4=11
\end{cases*}$の解は$x_k=$\fbox{\makebox[1.0cm][c]{$k$}}(k=1,2,3,4)である.
\begin{flushright}
    （2025線形代数学1a）
\end{flushright}
\fbox{2}
同次形の連立1次方程式を解け.\\
\[\begin{pmatrix}
    1 & -2 & 0 & 0 & 1 & 2 \\
    1 & -1 & 1 & 0 & 1 & 0 \\
    2 & -5 & -1 & 0 & 3 & 3 
\end{pmatrix}
\begin{pmatrix}
    x_1 \\ x_2 \\ x_3 \\ x_4 \\ x_5 \\ x_6 
\end{pmatrix}=
\begin{pmatrix}
    0 \\ 0 \\ 0
\end{pmatrix}\]
\begin{flushright}
    （三宅敏恒「線形代数概論」p62）
\end{flushright}
\fbox{3}
連立一次方程式
\[\begin{cases*}
    kx+(k^2-3k)z=3k-2 \\
    -2x-y+4z=-7 \hspace*{1.3cm}(k\text{は定数})\\
    x+y-3z=5
\end{cases*}\]
について,以下の空欄を適切な数式または数値で埋めよ.パタメーターを要する場合は,文字s,t,u,...をこの順に使用し,「標準的なパラメーターの置き方」に従うこと.
\begin{itemize}
    \item 解が存在しないならば$k=$\fbox{\makebox[1.0cm][c]{1}}である.
    \item 解が無数にあるならば$k=$\fbox{\makebox[1.0cm][c]{2}}で,このとき$x=$\fbox{\makebox[1.0cm][c]{3}}である.
    \item 解がただ一つならば$z=$\fbox{\makebox[1.0cm][c]{4}}である.
\end{itemize}
\begin{flushright}
    （2025線形代数学1a）
\end{flushright}
\fbox{4}
連立一次方程式が解を持つように,$a,b$を定めよ.\\
\[\begin{pmatrix}
    2 & 1 & 3 \\
    0 & -1 & 1 \\
    1 & 1 & 1
\end{pmatrix}
\begin{pmatrix}
    x_1 \\ x_2 \\ x_3
\end{pmatrix}=
\begin{pmatrix}
    1 \\ a \\ b
\end{pmatrix}\]
\begin{flushright}
    （三宅敏恒「線形代数概論」p64）
\end{flushright}
\subsubsection{解答}
\noindent
\fbox{1}\\(1)拡大係数行列より\\
$[A|\symbfit{b}]=
\begin{bmatrix}
    4 & 8 & 12 \\
    7 & 12 & 21
\end{bmatrix}\rightarrow
\begin{bmatrix}
    1 & 2 & 3 \\
    1 & 2 & 3 \\
\end{bmatrix}\rightarrow
\begin{bmatrix}
    1 & 2 & 3 \\
    0 & 0 & 0
\end{bmatrix} \hspace{1cm}\therefore
\begin{cases*}
    x+2s=3 \\
    y=s 
\end{cases*}\quad \rightarrow \quad
\begin{bmatrix}
    x \\ y
\end{bmatrix}=
s\begin{bmatrix}
    -2 \\ 1  
\end{bmatrix}+
\begin{bmatrix}
    3 \\ 0
\end{bmatrix}$\\\\
(2)拡大係数行列より\\
$[A|\symbfit{b}]=
\begin{bmatrix}
    3 & 4 & 1 & 10 & 15 \\
    2 & 3 & 0 & 7 & 11
\end{bmatrix} \rightarrow
\begin{bmatrix}
    1 & 1 & 1 & 3 & 4 \\
    2 & 3 & 0 & 7 & 11
\end{bmatrix} \rightarrow
\begin{bmatrix}
    1 & 1 & 1 & 3 & 4 \\
    0 & 1 & -2 & 1 & 3 
\end{bmatrix}\rightarrow
\begin{bmatrix}
    1 & 0 & 3 & 2 & 1 \\
    0 & 1 & -2 & 1 & 3 
\end{bmatrix}$\\\\
\hspace{2cm}
$\therefore \begin{cases*}
    x_1 +3s+2t=1 \\
    x_2 -2s+t=1
\end{cases*}\quad \rightarrow \quad
\symbfit{x}=s
\begin{bmatrix}
    -3 \\ 2 \\ 1 \\ 0
\end{bmatrix}+t
\begin{bmatrix}
    -2 \\ -2 \\ 0 \\ 1
\end{bmatrix}+
\begin{bmatrix}
    1 \\ 3 \\ 0 \\ 0
\end{bmatrix}$\\
コメント\\
ガウスの消去法を使うことを考えればREF行列までの変形で十分である.しかし,パラメーター付きの解は長く,代入する方が,計算がめんどくさくなる.
この程度の行列ならRREF行列まで変形させてしまうのも手である.\\

\noindent
\fbox{2}\\
拡大係数行列より\\
$[A|\symbfit{b}]=
\begin{bmatrix}
    1 & -2 & 0 & 0 & 1 & 2 \\
    1 & -1 & 1 & 0 & 1 & 0 \\
    2 & -5 & -1 & 0 & 3 & 3 
\end{bmatrix} \rightarrow
\begin{bmatrix}
    1 & -2 & 0 & 0 & 1 & 2 \\
    0 & 1 & 1 & 0 & 0 & -2 \\
    0 & -1 & -1 & 0 & 1 & -1 
\end{bmatrix}\rightarrow
\begin{bmatrix}
    1 & -2 & 0 & 0 & 1 & 2 \\
    0 & 1 & 1 & 0 & 0 & -2 \\
    0 & 0 & 0 & 0 & -1 & -3
\end{bmatrix}\\ \hspace{1cm}\rightarrow
\begin{bmatrix}
    1 & -2 & 0 & 0 & 0 & 1 \\
    0 & 1 & 1 & 0 & 0 & -2 \\
    0 & 0 & 0 & 0 & -1 & -3    
\end{bmatrix}\quad \therefore
\begin{cases*}
    x_1+2x_2+x_6=0\\
    x_2+x_3-2x_6=0 \\
    x_5-3x_6=0
\end{cases*}$\\
\qquad$\symbfit{x}=c_1
\begin{pmatrix}
    -2 \\ -1 \\ 1 \\ 0 \\ 0 \\ 0
\end{pmatrix}+c_2
\begin{pmatrix}
    0 \\ 0 \\ 0 \\ 1 \\ 0 \\ 0
\end{pmatrix}+c_3
\begin{pmatrix}
    -1 \\ 2 \\ 0 \\ 0 \\ 3 \\ 1
\end{pmatrix} \quad (c_1,c_2,c_3 \in \mathbf{R})$\\
\noindent
\fbox{3}\\
拡大係数行列より\\
$[A|\symbfit{b}]=
\begin{bmatrix}
    k & 0 & k^2-3k & 3k-2 \\
    -2 & -1 & 4 & -7 \\
    1 & 1& -3 & 5
\end{bmatrix}\rightarrow
\begin{bmatrix}
    k & 0 & k^2-3k & 3k-2 \\
    0 & 1 & -2 & 3 \\
    1 & 1 & -3 & 5
\end{bmatrix}\rightarrow
\begin{bmatrix}
    k & 0 & k^2-3k & 3k-2 \\
    0 & 1 & -2 & 3 \\
    1 & 0 & -1 & 2
\end{bmatrix}\\ \hspace{3cm}\rightarrow
\begin{bmatrix}
    1 & 0 & -1 & 2 \\
    0 & 1 & -2 & 3 \\
    0 & 0 & k(k-2) & k-2
\end{bmatrix}$\\\\
(1)解が存在しない時
$\begin{bmatrix}
    0 & 0 & 0 & \text{実数}(\neq 0)
\end{bmatrix}$となる.つまり$k=0$\\
このとき
$\begin{bmatrix}
    1 & 0 & -1 & 2 \\
    0 & 1 & -2 & 3 \\
    0 & 0 & 0 & 2
\end{bmatrix}$\\
(2)解が無数にある時rank($[A|\symbfit{b}]) < $列数の時であるから第3行がすべて0となる$k$を考える.つまり$k=2$\\
このとき
$\begin{bmatrix}
    1 & 0 & -1 & 2 \\
    0 & 1 & -2 & 3 \\
    0 & 0 & 0 & 0
\end{bmatrix}$より$x=s+2$\\
(3)線形方程式の解は3パターンで解なし,無数解,唯一解である.前2つはk=0,2の時なので,$k \neq 0,2$を考えるのが妥当である.\\
唯一解の時3行3列が1となればいいので,第3行を$k(k-2)$で割ると,
$\begin{bmatrix}
    1 & 0 & -1 & 2 \\
    0 & 1 & -2 & 3 \\
    0 & 0 & 1 & \frac{1}{k}
\end{bmatrix}$.
後退代入より$\displaystyle z=\frac{1}{k}$\\\\
コメント\\
rankが絡む問題はREF行列を活用すべき.REF行列から解を求める時はもちろん後退代入.\\なぜならREF行列への変形が前進消去なのだから.\\
\fbox{4}\\
拡大係数行列より\\
$[A|\symbfit{b}]=
\begin{bmatrix}
    2 & 1 & 3 & 1\\
    0 & -1 & 1 & a\\
    1 & 1 & 1 & b
\end{bmatrix}\rightarrow
\begin{bmatrix}
    0 & -1 & 1 & 1-2b \\
    0 & -1 & 1 & a \\
    1 & 1 & 1 & b
\end{bmatrix} \rightarrow
\begin{bmatrix}
    1 & 1 & 1 & b \\
    0 & -1 & 1 & a \\
    0 & 0 & 0 & 1-2b-a
\end{bmatrix}$\\
解を持つので$1-2b-a=0$
\subsection{行列式}
行列式の計算は知っていると優位になれるような行列や初見問題が出題されるようだ...
\subsubsection{問題}
\noindent
\fbox{1}\\
行操作を用いて,3\times 3の「ヴァンデルモンド(ヴァンデルモンデ)行列式」が\\
\[\det
\begin{bmatrix}
    1 & a & a^2 \\
    1 & b & b^2 \\
    1 & c & c^2
\end{bmatrix}=(b-a)(c-a)(c-b)\]
であることを示せ.
\begin{flushright}
    (ストラング：線形代数イントロダクションp268)
\end{flushright}
\fbox{2}\\
正定値対称行列Pは$P=LDL^T$と対角化することが可能である.また$A=BC$のとき$\det(A)=\det(B)\cdot\det(C)$が成り立つ.
次の正定値対称行列Pを定義する.\\
\[
P=
\begin{bmatrix}
    1 & 1 & 1 \\
    1 & 2 & 3 \\
    1 & 3 & 6
\end{bmatrix}\rightarrow
L=
\begin{bmatrix}
    1 & 0 & 0 \\
    1 & 1 & 0 \\
    1 & 2 & 1
\end{bmatrix}\quad D=
\begin{bmatrix}
    1 & 0 & 0 \\
    0 & 1 & 0 \\
    0 & 0 & 2
\end{bmatrix}からなる行列で対角化できる.\]
$\det(P)$を求めよ.
\begin{flushright}
    (自作)
\end{flushright}
\fbox{3}\\
AをUに簡約化して,$\det A$を求めよ.\\
\[A_1=
\begin{bmatrix}
    1 & 1 & 1 \\
    1 & 2 & 2 \\
    1 & 2 & 3
\end{bmatrix}\hspace{1cm}A_2=
\begin{bmatrix}
    1 & 2 & 3 \\
    2 & 2 & 3 \\
    3 & 3 & 3
\end{bmatrix}\]
\begin{flushright}
    (ストラング：線形代数イントロダクションp268より改題)
\end{flushright}
\subsubsection{解答}
\noindent
\fbox{1}\\
一般式は見づらい.ヴァンデルモンド行列は特徴的な形をしている.構成する要素の総積が答えである.\\
また総積とはとりうる組み合わせの差同士の積である.\\
一般式を付す.\\
\[V \coloneq
\begin{bmatrix}
    1 & x_1 & x_1^2 & \cdots & x_1^{n-1} \\
    1 & x_2 & x_2^2 & \cdots & x_2^{n-1} \\
    \vdots & \vdots & \vdots & \ddots & \vdots \\
    1 & x_n & x_n^2 & \cdots & x_n^{n-1} \\
\end{bmatrix}のとき,\det(V)=\det(V^T)=\prod_{i \leq i < j \leq n}^{}(x_j-x_i)\]
常に「下の行の要素ー上の要素」であることに注意.\\
\fbox{2}\\
サラスの公式より三角行列は0を通らない対角要素のみが積算される.また対角行列Dは対角要素のみが積算になる.\\
\[
P=LDL^T \Leftrightarrow 
\begin{bmatrix}
    1 & 1 & 1 \\
    1 & 2 & 3 \\
    1 & 3 & 6
\end{bmatrix} =
\begin{bmatrix}1 & 0 & 0 \\1 & 1 & 0 \\ 1 & 2 & 1\end{bmatrix} 
\begin{bmatrix}1 & 0 & 0 \\0 & 1 & 0 \\0 & 0 & 2\end{bmatrix}
\begin{bmatrix}1 & 1 & 1 \\0 & 1 & 2 \\0 & 0 & 1\end{bmatrix}
\]
よって$\det(L)=\det(L^T)=1$, $\det(D)=2$であるから,
\[\det(P)=\det(L)\cdot\det(D)\cdot\det(L^T)=1\cdot2\cdot1=2.\]
線形代数学2bでは「対角化」について学ぶ.対角化がいかに計算に対して嬉しい前処理かがわかる一端が見える.\\
\fbox{3}\\
(発展)で紹介したpivot積の話である.行の変換や因数を外に出すときの計算に注視してほしい.\\
\[\det(A_1) = \det
\begin{bmatrix}
    1 & 1 & 1 \\
    0 & 1 & 1 \\
    0 & 1 & 2
\end{bmatrix}= \det
\begin{bmatrix}
    1 & 0 & 0 \\
    0 & 1 & 1 \\
    0 & 0 & 1
\end{bmatrix}= \det
(I_3) \qquad\therefore \det(A_1) = 1\]\\
\hspace{2cm}$\det(A_2) = \det
\begin{bmatrix}
    1 & 2 & 3 \\
    0 & -2 & -3 \\
    0 & -3 & -6
\end{bmatrix}= (-1) \cdot (-3)\det
\begin{bmatrix}
    1 & 0 & 0 \\
    0 & 2 & 3 \\
    0 & 1 & 2
\end{bmatrix}= 3 \cdot \det
\begin{bmatrix}
    1 & 0 & 0 \\
    0 & 0 & -1 \\
    0 & 1 & 2
\end{bmatrix}$\\\\
ここで行の交換を行う.\\
\hspace{3.5cm}$=-3\cdot \det
\begin{bmatrix}
    1 & 0 & 0 \\
    0 & 1 & 2 \\
    0 & 0 & -1 
\end{bmatrix}=-3\cdot \det
\begin{bmatrix}
    1 & 0 & 0 \\
    0 & 1 & 0 \\
    0 & 0 & -1
\end{bmatrix} \qquad \therefore \det(A_2)= 3.$\\
サラスの公式を用いて良いが,シンプルな行列は簡約化をやって解くというのも1つである.\\
融合問題でrankと解と行列式が1つ行列に対して聞かれた場合の優位性は高い.簡約化する.この時に行列式ように,約分した数や
行の交換による-1などを$\rightarrow$の上とかにメモしておく.こうすればrankを求めた後に,素早く行列式を求めることが可能である.\\
またこの手法はコレスキー分解やLU分解でも使える.(コレスキー分解は「数値解析」で知れる)\\
※関係ないがコレスキー分解は正定値対称行列でのみ可能である.「数値解析」の課題で出される行列に注目.\\
※さらに正定値対称行列の条件も詳しく書いとこうかなと思ったけど踏みとどまる.\\
\subsection{線形独立と線形従属}
\noindent
\textbf{問題}\\
\fbox{1}

$\symbfit{v}_1,\;\symbfit{v}_2,\,\symbfit{v}_3$は線形独立であり,$\symbfit{v}_1,\,\symbfit{v}_2,\,\symbfit{v}_3,\,\symbfit{v}_4$は戦形従属であること示せ：
    \[
    \symbfit{v}_1 = \begin{bmatrix}1\\0\\1\end{bmatrix}, \qquad
    \symbfit{v}_2 = \begin{bmatrix}1\\1\\0\end{bmatrix}, \qquad
    \symbfit{v}_3 = \begin{bmatrix}1\\1\\1\end{bmatrix}, \qquad
    \symbfit{v}_4 = \begin{bmatrix}2\\3\\4\end{bmatrix}
    \]
\begin{flushright}
    (ストラング：線形代数イントロダクション p.189)
\end{flushright}
\noindent
\fbox{2}

以下のうち,線形独立なベクトルの個数の最大値を求めよ.
\[
    \symbfit{v}_1 = \begin{bmatrix}1\\-1\\0\\0 \end{bmatrix}, \quad 
    \symbfit{v}_2 = \begin{bmatrix}1\\0\\-1\\0 \end{bmatrix}, \quad 
    \symbfit{v}_3 = \begin{bmatrix}1\\0\\0\\1 \end{bmatrix}, \quad 
    \symbfit{v}_4 = \begin{bmatrix}0\\1\\-1\\0 \end{bmatrix}, \quad 
    \symbfit{v}_5 = \begin{bmatrix}0\\1\\0\\-1 \end{bmatrix}, \quad 
    \symbfit{v}_6 = \begin{bmatrix}0\\0\\1\\-1 \end{bmatrix} 
\]
\begin{flushright}
    (ストラング：線形代数イントロダクション p.189)
\end{flushright}
\fbox{3}

$
    \symbfit{a}_1 = \begin{bmatrix}1\\-1\\2 \end{bmatrix}, \quad 
    \symbfit{a}_2 = \begin{bmatrix}-1\\4\\1 \end{bmatrix}, \quad 
    \symbfit{a}_3 = \begin{bmatrix}3\\-4\\5 \end{bmatrix}
$
が線形従属であることを示し,$\symbfit{a}_2$を$\symbfit{a}_1$と$\symbfit{a}_3$の線形結合で示せ.
\begin{flushright}
    (マセマ：線形代数キャンパス・ゼミ p.139)
\end{flushright}
\fbox{4}

$
    \symbfit{u}_1 = \begin{bmatrix}1\\1\\0\\1 \end{bmatrix}, \quad 
    \symbfit{u}_2 = \begin{bmatrix}0\\1\\2\\1 \end{bmatrix}, \quad 
    \symbfit{u}_3 = \begin{bmatrix}1\\-1\\1\\0 \end{bmatrix}, \quad
    \symbfit{u}_4 = \begin{bmatrix}2\\0\\-1\\1 \end{bmatrix}
$が,$\mathbf{R}^4$(4次元空間)において,1組の基底となることを示せ.
\begin{flushright}
    (マセマ：線形代数キャンパス・ゼミ p.140)
\end{flushright}

\noindent
\textbf{解答}\\
\fbox{1}

$A=\begin{bmatrix}
    && \\
    \;\symbfit{v}_1 & \symbfit{v}_2 & \symbfit{v}_3 \\
    && 
\end{bmatrix}$において$\det(A)\neq 0$より示せた.また,$B=\begin{bmatrix}
    &&& \\
    \;\symbfit{v}_1 & \symbfit{v}_2 & \symbfit{v}_3 & \symbfit{v}_4 \\
    &&& 
\end{bmatrix}$において行列Bのサイズは$3\times 4$より線形従属となる.\\
\fbox{2}

4行なので線形独立の列ベクトルは理論上4つは超えない.-1の位置から考えて$\symbfit{v}_1,\;\symbfit{v}_2,\,\symbfit{v}_3$が上がる.
他のベクトルは今あげた3つの列ベクトルで表現できることが容易な計算でわかるので,簡約化せずに3つとなる.\\
\fbox{3}

\[
A=\begin{bmatrix}1 & -1 & 3 \\ -1 & 4 & -4 \\2 & 1 & 5\end{bmatrix}
\quad \rightarrow \quad
U=\begin{bmatrix}1 & -1 & 3 \\ 0 & 3 & -1 \\ 0 & 0 & 0\end{bmatrix}\]
行列$A$は3列に対して$rank(A)=2$より線形従属となる.また,
$U\symbfit{c}=\mathbf{0}$より
\[
\begin{cases*}
    c_1 -c_2+3c_3 = 0 \\
    3c_2-c_3=0
\end{cases*} \qquad \therefore \;a_2=8a_1-3a_3.\]
\fbox{4}

\[
A=\begin{bmatrix}
    1 & 0 & 1 & 2 \\ 1 & 1 & =1 & 0 \\ 0 & 2 & 1 & -1 \\ 1 & 1 & 0 & 1\end{bmatrix}
    \quad \rightarrow \quad    
B=\begin{bmatrix}
    1 & 0 & 1 & 2 \\ 0 & 1 & -2 & -2 \\ 0 & 0 & 1 & 1 \\ 0 & 0 & 0 & -2
\end{bmatrix}\]
つまり非退化より線形独立であるため,基底となる.
\subsection{固有値と固有ベクトル}
\noindent
\textbf{問題}\\
\fbox{1}

以下の$A$と$A^{-1}$の固有値と固有ベクトルを計算せよ.さらに,発展を読んだ人はトレースを確かめよ.
\[
A=\begin{bmatrix}0 & 2 \\ 1 & 1\end{bmatrix}
\quad\text{と}\quad
A^{-1}=\begin{bmatrix}-\frac{1}{2} & 1 \\ \frac{1}{2} & 0\end{bmatrix}
\]
\begin{flushright}
    (ストラング：線形代数イントロダクション p.312)
\end{flushright}
\noindent
\textbf{解答}\\
\fbox{1}
\[
g_A(\lambda)=\begin{vmatrix}
    0-\lambda & 2 \\ 1 & 1-\lambda
\end{vmatrix}=\lambda(\lambda-1)-2=(\lambda-2)(\lambda+1)\qquad \therefore\;\lambda=-1,\;2 \;(g_A(\lambda)=0)\]
$\lambda=-1$のとき
\[
\begin{bmatrix}1 & 2 \\ 1 & 2\end{bmatrix} 
\quad \rightarrow \quad
\begin{bmatrix}1 & 2 \\ 0 & 0\end{bmatrix} \qquad 
\text{固有ベクトルは}\begin{bmatrix}
    -2 \\ 1
\end{bmatrix}
\]
$\lambda=2$のとき
\[
\begin{bmatrix}-2 & 2 \\ 1 & -1\end{bmatrix} 
\quad \rightarrow \quad
\begin{bmatrix}1 & -1 \\ 0 & 0\end{bmatrix} \qquad 
\text{固有ベクトルは}\begin{bmatrix}
    1 \\ 1
\end{bmatrix}
\]
\[
g_{A^{-1}}(\lambda)=\begin{vmatrix}
    -\frac{1}{2}-\lambda & 1 \\ \frac{1}{2} & 0-\lambda
\end{vmatrix}=\lambda(\lambda+\frac{1}{2})-\frac{1}{2}=(\lambda-\frac{1}{2})(\lambda+1)\qquad \therefore\;\lambda=-1,\;\frac{1}{2} \;(g_{A^{-1}}(\lambda)=0)\]
$\lambda=-1$のとき
\[
\begin{bmatrix}\frac{1}{2} & 1 \\ \frac{1}{2} & 1\end{bmatrix}
\quad \rightarrow \quad
\begin{bmatrix}-\frac{1}{2} & 1 \\ 0 & 0\end{bmatrix} \qquad
\text{固有ベクトルは}\begin{bmatrix}
    -2 \\ 1
\end{bmatrix}
\]
$\lambda=\frac{1}{2}$のとき
\[
\begin{bmatrix}-1 & 1 \\ \frac{1}{2} & -\frac{1}{2}\end{bmatrix} 
\quad \rightarrow \quad
\begin{bmatrix}-1 & 1 \\ 0 & 0\end{bmatrix} \qquad 
\text{固有ベクトルは}\begin{bmatrix}
    1 \\ 1
\end{bmatrix}
\]
\begin{tcolorbox}[colframe=black!80,colback=white,title=摘要]
    \hspace{1em}
    対角化で逆行列を素早く計算できる話が出た.これを見ると逆行列の固有値は元の行列の固有値の逆数になることや,
    固有ベクトルが変わらないことはすぐわかる.なので,テキスト範囲では答えが「予想通り」という問題ではある.
    ただ,一般化の証明は難しいところがある.
\end{tcolorbox}
\subsection{修正グラム--シュミットについて}
\noindent
\textbf{問題}\\
\hspace{1em}
以下のベクトルの組にグラム=シュミットの正規直行化を適用せよ.問題文のベクトルの順序を変えないこと.分母の有理化は不要.\\
    \[\begin{pmatrix}1 \\ 0 \\ 0 \\ 1\end{pmatrix},\;
    \begin{pmatrix}1 \\ 1 \\ 1 \\ 0\end{pmatrix},\;
    \begin{pmatrix}1 \\ 0 \\ 1 \\ 0\end{pmatrix}\]
\begin{flushright}
    (2024年 線形代数学2a期末試験)
\end{flushright}
\vspace{2cm}
\textbf{解答}\\
\hspace{1em}
3つのベクトルを順に$\symbfit{a},\symbfit{b},\symbfit{c}$とする.
\paragraph{1. $\symbfit{q}_1$ の計算}
    \[
    \symbfit{a} = \symbfit{A} \quad \rightarrow \quad
    \symbfit{q}_1 = \frac{\symbfit{A}}{\|\symbfit{A}\|}
        = \frac{1}{\sqrt{1^2+0^2+0^2+1^2}}\,\symbfit{a}
        = \frac{1}{\sqrt{2}}\begin{pmatrix}1 \\ 0 \\0 \\ 1\end{pmatrix}.
    \]

    \paragraph{2. $\symbfit{q}_2$ の計算}
    \[
        \symbfit{B} = \symbfit{b} - (\symbfit{q}_1^T \symbfit{b})\,\symbfit{q}_1,
        \qquad
        \symbfit{q}_1^T \symbfit{b} = \frac{1}{\sqrt2}(1+0+0+0)=\frac{1}{\sqrt2},
    \]
    \[
    \symbfit{B}
        = \begin{pmatrix}1 \\ 1 \\ 1 \\ 0\end{pmatrix}
         - \frac{1}{\sqrt2}\cdot\frac{1}{\sqrt2}
        \begin{pmatrix}1 \\ 0 \\ 0 \\ 1\end{pmatrix}
        =\frac{1}{2}\left\{
        \begin{pmatrix} 2 \\ 2 \\ 2 \\ 0 \end{pmatrix}
        - \begin{pmatrix} 1 \\ 0 \\ 0 \\ 1 \end{pmatrix}\right\}
        =\frac{1}{2}
        \begin{pmatrix} 1 \\ 2 \\ 2 \\ -1 \end{pmatrix},
    \]
    \[
        \symbfit{q}_2 = \frac{\symbfit{B}}{\|\symbfit{B}\|}
        = \frac{\cancel{\frac{1}{2}}}{\cancel{\frac{1}{2}}\sqrt{1^2+2^2+2^2+(-1)^2}}
        \begin{pmatrix}1 \\ 2 \\ 2 \\ -1\end{pmatrix}
        = \frac{1}{\sqrt{10}}
        \begin{pmatrix}1 \\ 2 \\ 2 \\ -1\end{pmatrix}.
    \]

    \paragraph{3. $\symbfit{q}_3$ の計算}
    \[
        \symbfit{C} = \symbfit{c} - (\symbfit{q}_1^T \symbfit{c})\,\symbfit{q}_1 - (\symbfit{q}_2^T \symbfit{c})\,\symbfit{q}_2,
    \]
    \[
        \symbfit{q}_1^T \symbfit{c} = \frac{1}{\sqrt2}(1+0+0+0)=\frac{1}{\sqrt{2}},\qquad
        \symbfit{q}_2^T \symbfit{c}
        = \frac{1}{\sqrt{10}}\begin{pmatrix} 1 \\ 2 \\ 2 \\ -1 \end{pmatrix}^T \begin{pmatrix} 1 \\ 0 \\ 1 \\ 0 \end{pmatrix}
        = \frac{3}{\sqrt{10}}
    \]
    \[
        \symbfit{C}
        = \begin{pmatrix} 1 \\ 0 \\ 1 \\ 0\end{pmatrix}
        - \frac{1}{\sqrt{2}} \cdot\frac{1}{\sqrt{2}}\begin{pmatrix} 1 \\ 0 \\ 0 \\ 1 \end{pmatrix}
        - \frac{3}{\sqrt{10}} \cdot\frac{1}{\sqrt{10}} \begin{pmatrix} 1 \\ 2 \\ 2 \\ -1 \end{pmatrix}
        =\frac{1}{10}\left\{
            \begin{pmatrix} 10 \\ 0 \\ 10 \\ 0 \end{pmatrix}
            - \begin{pmatrix} 5 \\ 0 \\ 0 \\ 5 \end{pmatrix}
            - \begin{pmatrix} 3 \\ 6 \\ 6 \\ -3 \end{pmatrix}
        \right\}
        =\frac{1}{5}\begin{pmatrix}
            1 \\ -3 \\ 2 \\ -1
        \end{pmatrix}
    \]
    \[
        \symbfit{q}_3
        = \frac{\symbfit{C}}{\|\symbfit{C}\|}
        = \frac{\cancel{\frac{1}{5}}}{\cancel{\frac{1}{5}}\sqrt{1^2+(-3)^2+2^2+(-1)^2}}\begin{pmatrix} 1 \\ -3 \\ 2 \\ -1 \end{pmatrix}
        = \frac{1}{\sqrt{15}}\begin{pmatrix} 1 \\ -3 \\ 2 \\ -1 \end{pmatrix}.
    \]
\subsection{対角化と冪乗}
\noindent
\textbf{問題}\\
実対称行列
$
A=
\begin{pmatrix}
  6 & 3 \\
  3 & -2
\end{pmatrix}$
を直交行列で対角化し,$A^n (n = 1,2, \dots )$を計算せよ.以下の指示に従い,[1]-[10]を埋める形で解答すること.\\
【[1][2]】は左から小さい値.【[3][4]】整数部分のベクトルを.更に,第一成分は非負(ゼロまたは正の数)とすること
【[5][6]】基底に対して正規直交化を行い,ベクトル成分は整数とすること.例えば,
$\begin{pmatrix}
  1/\sqrt{3} \\
  1/\sqrt{3} \\
\end{pmatrix}$
を記入したいなら
$\displaystyle
\frac{1}{\sqrt{3}}
\begin{pmatrix}
  1 \\
  2
\end{pmatrix}$
と記すこと.
\[
\begin{array}{c|c|c}
  \text{固有値} & \quad  \text{[1]} \quad & \quad \text{[2]} \quad \\
  \hline
  \text{固有ベクトル} & \quad \text{[3]} \quad & \quad \text{[4]} \quad \\
  \hline
  \text{直交化したベクトル} & \quad \text{[5]} \quad & \quad \text{[6]} \quad \\
\end{array}
  \qquad P = \frac{1}{\sqrt{10}}
  \begin{pmatrix}
    & & \\
    \quad & [7] & \quad \\
    & & 
  \end{pmatrix}
   \text{ならば} P^{-1}AP = 
   \begin{pmatrix}
    & & \\
    \quad & [8] & \quad \\
    & & 
  \end{pmatrix}\\
\]
\[PD^nP^{-1}=
\frac{[1]^n}{10}
\begin{pmatrix}
  & & \\
  \quad & [9] & \quad \\
  & & 
\end{pmatrix}
+
\frac{[2]^n}{10}
\begin{pmatrix}
  & & \\
  \quad & [10] & \quad \\
  & & 
\end{pmatrix}
\]
\begin{flushright}
    (2024年 線形代数学2b期末試験)
\end{flushright}

\noindent
\textbf{解答}\\
\hspace{1em}
固有多項式$g_A(t)$より
\[ 
g_A(t)=
\begin{vmatrix}
  t-6 & -3 \\
   -3 & t+2 
\end{vmatrix}
=(t-7)(t+3) \qquad \therefore t = -3,7
\]
t = $-3$ のとき
\[
-3E - A = 
\begin{pmatrix}
  -9 & -3 \\
  -3 & -1
\end{pmatrix}\]
\[
\begin{array}{cc} 
  3 & 1 \\    
  3 & 2 \\
  \hline
  3 & 1 \\
  0 & 0
\end{array}
\qquad \qquad 
\begin{cases} 
  3 x + \alpha = 0\\
  y = \alpha
\end{cases}
\qquad
\therefore
\alpha
\begin{pmatrix}
  - \frac{1}{3}\\
  1 \\
\end{pmatrix}
\rightarrow \alpha
\begin{pmatrix}
  1\\
  -3 
\end{pmatrix}
\] 
t = 7 のとき
\[
7E - A = 
\begin{pmatrix}
  1 & -3 \\
  -3 & 9
\end{pmatrix}\]
\[
\begin{array}{cc}
  1 & -3 \\    
  1 & -3 \\
  \hline
  1 & -3 \\
  0 & 0
\end{array}
\qquad \qquad
\begin{cases}
  x - 3 \alpha = 0\\
  y = \alpha
\end{cases}
\qquad
\therefore
\alpha
\begin{pmatrix}
  3\\
  1 \\
\end{pmatrix}
\] 
\[
\begin{array}{c|c}
  \text{e-val} & \text{e-vec} \\ 
  \hline
  -3^{[1]} & 
  \begin{pmatrix}
    1\\
    -3
\end{pmatrix}^{[3]} 
\rightarrow \displaystyle
\frac{1}{\sqrt{10}}
\begin{pmatrix}
  1\\
  -3
\end{pmatrix}^{[5]} \\
\hline
 7^{[2]} & 
 \begin{pmatrix}
  3\\
  1
 \end{pmatrix}^{[4]}
 \rightarrow \displaystyle
\frac{1}{\sqrt{10}}
\begin{pmatrix}
  3\\
  1
\end{pmatrix}^{[6]} \\
\end{array}
\]
\begin{align*} 
  P = 
  \frac{1}{\sqrt{10}}
  \begin{pmatrix}
    1 & 3  \\
    -3 & 1 \\
  \end{pmatrix}^{[7]}
  \text{ならば}
  P^{-1}AP = 
  \begin{pmatrix}
    -3 &  \\
     & 7 \\
  \end{pmatrix}^{[8]}
\end{align*}
\begin{align*}
  \begin{pmatrix}
    1\\
    -3
  \end{pmatrix}
  \cdot
  \begin{pmatrix}
    1 & -3
  \end{pmatrix}
  = 
  \begin{pmatrix}
    1 & -3 \\
    -3 & 9
  \end{pmatrix}^{[9]}
\end{align*}
\begin{align*}
  \begin{pmatrix}
    3\\
    1
  \end{pmatrix}
  \cdot
  \begin{pmatrix}
    3 & 1
  \end{pmatrix}
  = 
  \begin{pmatrix}
    9 & 3 \\
    3 & 1
  \end{pmatrix}^{[10]}
\end{align*}
  \hspace{1em}
\begin{tcolorbox}[colframe=black!80,colback=white,title=適要]
\[
A=PDP^{-1}
\]
また,実対称行列において$P^{-1}=P^T$が成り立つ.そのため,$P^TP=PP^T=E$が成り立つ.(このとき,$A^T$はAの転置行列を意味する)
\paragraph{\;\small $P^{-1}APについて$}
\begin{align*}
  A &= PDP^{-1} \\
  AP & = PDP^{-1}P = PDE = PD \quad (右からP\;をかけた) \\
  P^{-1}AP &= P^{-1}PD \;\Leftrightarrow P^{-1}AP = D \quad (左からP^{-1}をかけた)
\end{align*}
\paragraph{\;\small $PD^nP^{-1}について$}
\begin{align*}
  A^n & = PDP^{-1} \cdot PDP^{-1}\cdot \; \cdots \; \cdot PDP^{-1} \\
      &(P^{-1}P=Eより)\\
      & = PDED \dots P^{-1} \\
      & = PD^nP^{-1}
\end{align*}
\end{tcolorbox}
\subsection{サービス問題}
※この問題はかなり難しいかもしれない.読む価値はあると思うので,すぐ答えを見てもいい.
\vspace{5mm}\\
\hspace{1em}
とある多変数関数$f(x)$と関数上の点$u$について議論する.(x,uは列ベクトル)\\
$f(x)$を$u$についてのテイラー展開で2次まで近似すると,
\[
f(x)=f(u)+(u-x)^T \frac{\partial f}{\partial x}\;u+\frac{1}{2}(u-x)^T\;H(u)(u-x)\]
と表現できる.このとき,$H(u)$は2回微分が含まれる.\\
2回微分は曲率を意味し,「どれくらい曲がっているか」を表す.そのため,Hの要素は各方向の曲率が入る.\\
つまり,点$u$が変曲点の場合を考えている.では$\displaystyle\frac{\partial}{\partial x}\;f(u)=0$であるとき,$\displaystyle(u-x)^T \frac{\partial f}{\partial x}\;u$はゼロとなるため,$\displaystyle\frac{1}{2}(u-x)^T\;H(u)(u-x)$が決定項となる.\\
H(u)が正定値(全方向に対して凸)であるとき,$f(u)$は$f(x)$にとってどんな意味を持つか.\\
※正の曲がりを凸,負の曲がりを凹という.
\begin{flushright}
    (ストラング：計算理工学p.66-68から自作)
\end{flushright}
\vspace{6.52mm}
\textbf{解答}\\
\hspace{1em}
単刀直入にいうと,\textbf{極小点}となる.$\frac{\partial f}{\partial x}\;u$は1回微分なので関数の傾きを表す.これがゼロなので,極値もしくは鞍点である.
この2択を判別するのに使うのが行列Hである.Hが正定値であるとき,全方向に対して凸を意味し,$f(x)=\frac{1}{2}(u-x)^T\;H(u)(u-x)$より$x$が$u$に近づくほど$f(x)$は最小になる.つまり,点$u$のとき極小点になるということである.\\
鞍点は見る角度によっては極小値(極大値)という点なので,方向によっては凸ではない.つまり,Hが半正定値(一部方向が凸もしくはゼロ)あるいは不定値(一部方向が凸もしくは凹)のときはその点は鞍点となる.このとき,$\frac{1}{2}(u-x)^T\;H(u)(u-x)$の正負判定ができない.

ちなみにこの行列を\textbf{ヘッシアン}$_{\text{Hessian}}$という.
\[
H_{i\,j}=\frac{\partial^2 f}{\partial u_i\;\partial u_j}=\frac{\partial^2 f}{\partial u_j\;\partial u_i}=H_{j\;i}\]
微積の授業ではヘッシアンの行列式を計算したと思う.これはまさしく正定値かどうかの判別をしていたのである.\\
以下に極小値もしくは鞍点のときの$u^THu$のグラフの概形を示す.Kは正定値,Bは半正定値,Mは不定値である.
\begin{center}
    \includegraphics*[width=0.8\textwidth]{Hessian.png}
\end{center}
まとめると,
\begin{itemize}
    \item $\frac{\partial f}{\partial x}\;u=0$かつ,Hが正定値のとき,$(x-u)^T H(u)(x-u)>0$より$x=u$のとき極小値
    \item $\frac{\partial f}{\partial x}\;u=0$かつ,Hが正定値でないとき,$(x-u)^T H(u)(x-u)$の正負がわからないので$x=u$のとき傾きは0だが,極小値ではない.つまり鞍点.
\end{itemize}
※全方向に凹である場合,Hは負定値となる.このときは極大値となる.
